\documentclass[conference]{IEEEtran}
\usepackage{amsmath,amssymb,graphicx,url,array}
\raggedbottom
\makeatletter\renewenvironment{table}[1][]{\par\noindent\begin{minipage}{\linewidth}\def\@captype{table}\footnotesize}{\end{minipage}\par}\makeatother
\begin{document}
\title{Literature Review: Speculative Decoding for Lossless Language Model Acceleration}
\author{\IEEEauthorblockN{Anonymous}}
\maketitle
\begin{abstract}
Speculative sampling accelerates autoregressive LLM decoding by generating candidate tokens from a draft model and accepting them via modified rejection sampling against a larger target model, achieving speedups up to 2.5× across benchmarks. \cite{chen2023accelerating} Speculative decoding accelerates transformer inference by generating multiple tokens via a lightweight approximation model without altering output distribution. \cite{leviathan2022fast}
\end{abstract}
\begin{IEEEkeywords}speculative decoding, lossless acceleration, large language model, decoding efficiency, transformer inference\end{IEEEkeywords}
\section{Introduction}
Chen et al. proposed speculative sampling for accelerating LLM decoding, using a draft model to generate multiple tokens per target model call, with modified rejection sampling preserving target distribution within numerics [chen2023accelerating, p.1]. \cite{chen2023accelerating}

Yaniv Leviathan, Matan Kalman, and Yossi Matias propose speculative decoding to accelerate inference from large autoregressive Transformers. \cite{leviathan2022fast}

\section{Reviewed sources}
\subsection{Accelerating Large Language Model Decoding with Speculative Sampling}
Speculative sampling accelerates LLM decoding using a faster draft model to generate candidate tokens, which a larger target model scores and accepts via modified rejection sampling. Reported speedups of 2-2.5× on Chinchilla 70B and up to 2.46× on HumanEval with nucleus sampling, preserving target distribution without model modifications. \cite{chen2023accelerating}

\subsection{Fast Inference from Transformers via Speculative Decoding}
The source proposes speculative decoding using a small approximation model to generate candidate tokens in parallel, accelerating inference from large autoregressive Transformers without altering output distribution. Reported results show 2X–3X wall-time improvements over standard T5X on T5-XXL with identical outputs. \cite{leviathan2022fast}

\section{Source comparison}
\begin{table}[H]\centering
\begin{tabular}{|>{\raggedright\arraybackslash}p{0.26\linewidth}|>{\raggedright\arraybackslash}p{0.235\linewidth}|>{\raggedright\arraybackslash}p{0.235\linewidth}|}\hline
Source & Approach & Evaluation \\ \hline
Chen 2023 \cite{chen2023accelerating} & modified rejection & decoding speedup \\ \hline
Leviathan 2022 \cite{leviathan2022fast} & speculative decoding & wall-time comparison \\ \hline
\end{tabular}\caption{Qualitative comparison of the processed sources.}\end{table}
\section{Evidence limitations}
The source authors state: "Note that as \ensuremath{\mathit{K}}
increases, the overall speedup plateaus or even regresses, with XSum being optimal at \ensuremath{\mathit{K}} = 3." \cite{chen2023accelerating}

The source authors state: "One limitation of speculative execution in general, and of
speculative decoding in particular, is that latency is improved through increased concurrency at the cost of an increased number of arithmetic operations." \cite{leviathan2022fast}

\section{Conclusion}
Speculative sampling accelerates LLM decoding without modifying target model parameters, achieving significant speedups across benchmarks with appropriate draft models. \cite{chen2023accelerating}

The narrow takeaway is that speculative decoding provides modest but reliable speedups for large transformer inference when paired with a lightweight approximation model. \cite{leviathan2022fast}

\bibliographystyle{IEEEtran}
\bibliography{refs}\end{document}
